% section_4_additional.tex — Additional Analyses
% Drafted per: docs/Draft/THESIS_SKELETON.md v5 §IV (post-capex-drop 2026-04-22):
%   §IV.1 Drivers of CEO speech uncertainty (H11 PRisk + H24 US EPU + H24b GEPU) — reverse-direction construct validation
%   §IV.2 Outside-world reaction (H14c 25-day post-call closing-quote spread) — market info-asymmetry response
% Pending task: T43 (Draft §IV Additional Analyses)

\section{Additional Analyses}
\label{sec:additional}

\subsection{Drivers of CEO Speech Uncertainty}
\label{sec:additional:drivers}

The main analyses of \S\ref{sec:main} treat CEO speech uncertainty as the regressor. This subsection inverts the direction and asks whether the speech-uncertainty metric itself responds to exogenous uncertainty shocks, a construct-validation test for the main IV. If the metric captures economically meaningful uncertainty content rather than stylistic noise, it should covary positively with independently measured uncertainty at firm, US, and global levels of aggregation.

We estimate three reverse-direction regressions. The firm-level driver is PRisk, the political-risk count of \citeA{hassan2020}, which flags bigrams in a firm's earnings-call transcript that lie within ten words of a risk or uncertainty synonym and overlap a political-topic dictionary; PRisk therefore varies at the firm-quarter level, and the associated regressions cluster standard errors at the firm. The US macro driver is the news-based Economic Policy Uncertainty index of \citeA{baker2016}, matched to each call by calendar month; the global aggregation is the GDP-weighted Global EPU index of \citeA{davis2016}, likewise matched monthly. The macro regressions cluster standard errors two-way at the firm and calendar-quarter levels and absorb calendar-year fixed effects, so identification runs off within-year monthly variation in the driver. All three regressions are tested one-tailed in the theoretically predicted positive direction (exogenous uncertainty shocks raise measured speech uncertainty). Table~\ref{tab:driver_matrix} consolidates the three drivers against the two CEO speech-uncertainty channels (UncAnsCEO, UncPreCEO) under both industry and firm fixed effects.

The firm-level PRisk driver is strongly significant in every cell. Contemporaneous coefficients range from $+0.00013$ to $+0.00030$ per unit PRisk across the four CEO cells (Q\&A and Presentation channels $\times$ industry and firm fixed effects), with $t$-statistics between $9.0$ and $17.0$ and sample sizes of $65{,}394$ to $65{,}760$ firm-quarters. The PRisk-on-speech mapping therefore holds at high statistical precision across both CEO channels.

The macro drivers exhibit a channel-asymmetric pattern. US EPU raises UncPreCEO by $+0.021$ to $+0.024$ per unit log-EPU ($p$-values $0.0008$--$0.0021$, one-tailed); the Q\&A coefficient is smaller, $+0.012$ to $+0.014$, and reaches conventional significance only under firm fixed effects ($p=0.037$; the industry-FE cell is $p=0.052$). Global EPU follows the same pattern at larger magnitudes: UncPreCEO loads at $+0.027$ (industry FE, $p=0.0056$) to $+0.031$ (firm FE, $p=0.0021$); UncAnsCEO loads at $+0.018$ (industry FE, $p=0.051$) and $+0.021$ (firm FE, $p=0.027$).

The channel asymmetry here is inverted relative to the cash regressions of \S\ref{sec:main}: the Q\&A channel dominates in \S\ref{sec:main:hc} and \S\ref{sec:main:hfc}, but the Presentation channel dominates in the macro-driver regressions. The firm-level PRisk measure, which is itself constructed at the firm-call level, loads both channels approximately equally. The pattern is descriptively consistent with prepared remarks carrying broad uncertainty content (macroeconomic, political) and spontaneous Q\&A content responding more tightly to firm-specific conditions; a fuller interpretive treatment is deferred to \S\ref{sec:conclusion}. Comparable results for the NoCEO speech channels, which identify whether uncertainty content in non-CEO speech also responds to the three drivers, are reported in Appendix~\ref{app:additional:drivers}.

\subsection{Outside-World Reaction}
\label{sec:additional:reaction}

The preceding subsections show that CEO speech uncertainty co-moves with firm-level cash holdings (\S\ref{sec:main:hc}, \S\ref{sec:main:hfc}) and responds to exogenous uncertainty shocks (\S\ref{sec:additional:drivers}). The complementary question is whether agents external to the firm recognize CEO speech uncertainty as a signal and react to it. The canonical market channel is information asymmetry, which liquidity providers price into the bid-ask spread of the firm's stock: a firm whose CEO conveys more uncertainty in its earnings call should have wider subsequent bid-ask spreads as market-makers widen quotes against informed counterparties.

We test this through the daily closing-quote relative bid-ask spread averaged over the call day and the 25 following trading days, denoted Spread\textsubscript{25D}. The daily spread uses the standard closing-quote relative formulation, $2(\text{ASK}_d - \text{BID}_d)/(\text{ASK}_d + \text{BID}_d)$, computed from CRSP daily bid and ask quotes; the call-quarter observation is the simple mean across the call day and the 25 following trading days (a 26-day window), pooled-winsorized at the $1$st and $99$th percentiles. We adopt this window---wider than the 3-day announcement window typically used in short-horizon event studies---to capture market reaction sustained over approximately one trading month, beyond announcement-day volatility. The estimation uses Suite H14c, which enters the four call-level uncertainty IVs (UncAnsCEO, UncPreCEO, UncAnsMgr, UncPreMgr) jointly; we retain the manager-pool IVs here, unlike the CEO-only specification of \S\ref{sec:main}, because the post-call bid-ask spread is plausibly priced against both CEO-specific and broader manager-pool speech content, and partialling out the manager-pool channel isolates the CEO-specific contribution to the market's revision of information asymmetry. We report the CEO coefficients in the body and the manager-pool coefficients in Appendix~\ref{app:additional:reaction}. Standard errors are clustered at the firm. Inference is one-tailed positive on each speech-uncertainty IV, consistent with the information-asymmetry prediction.

The Presentation channel (UncPreCEO) raises the contemporaneous post-call spread. Across the six contemporaneous specifications (two fixed-effect families $\times$ three control sets) the coefficient ranges from $+0.04$ to $+0.43$; three of the six cells reach $p<0.10$ (one-tailed), two reach $p<0.05$, and one reaches $p<0.01$. The strongest cells are the base specifications under industry and firm fixed effects ($+0.30$ at $p=0.021$ and $+0.43$ at $p=0.009$ respectively). Sample sizes range from $44{,}092$ to $63{,}899$ firm-quarters across the six cells. On the one-quarter-lead dependent variable, the Presentation channel does not reach significance in any cell (coefficients $-0.03$ to $-0.24$, all negative); under the one-tailed positive convention, the Presentation-channel effect on the post-call spread is therefore confined to the contemporaneous call quarter.

The Q\&A channel (UncAnsCEO) does not move the contemporaneous spread: zero of six cells reach significance and point estimates span $-0.10$ to $+0.07$. It does, however, raise the one-quarter-lead spread across every specification. The lead coefficient ranges from $+0.18$ to $+0.44$; five of the six cells are significant at $p<0.10$ and two reach $p<0.05$. The one non-significant lead cell misses by less than half a percentage point ($p=0.100$) and is sign-aligned with the five significant cells. Sample sizes for the lead specifications range from $44{,}569$ to $63{,}443$ firm-quarters.

The two CEO channels therefore produce a temporally complementary pattern: prepared remarks widen the post-call spread within the same call quarter, while spontaneous Q\&A remarks widen the subsequent call-quarter spread. The contemporaneous Presentation effect is descriptively consistent with immediate price-formation reaction to the prepared statement of uncertainty, and the lagged Q\&A effect with slower market assimilation of analyst-question-conditional content; fuller interpretation is deferred to \S\ref{sec:conclusion}. The two channel asymmetries identified in this section---Presentation dominates macro-driver regressions in \S\ref{sec:additional:drivers} and contemporaneous spread here, while Q\&A dominates the cash-holdings regressions of \S\ref{sec:main:hc} and \S\ref{sec:main:hfc} and the lagged spread here---are descriptively consistent with the two segments of the call carrying distinct, economically meaningful uncertainty content that propagates through different downstream margins.
